\section{Introduction}
\label{sec:intro}


Weight decay is a canonical hyperparameter in deep learning whose role has evolved alongside changes in training regimes. In classical multi-epoch training, weight decay was understood primarily as a regularizer that improves generalization by shrinking weights and controlling model capacity \citep{hardt2016train, zhang2017understanding, sun2025investigating}. In contemporary large-scale pretraining, which often involves a single epoch over massive datasets \citep{brown2020language, kaplan2020scaling}, weight decay no longer primarily serves the purpose of generalization but plays a decisive role in optimization stability and convergence \citep{andriushchenko2023whywd, zhangdoes, wang2024set}. 

Moreover, modern large language models are typically developed in two distinct stages: a large-scale pretraining phase followed by a post-training phase involving supervised fine-tuning, alignment, and reinforcement learning \citep{brown2020language,instruct_gpt, bi2024deepseek,lambert2024tulu}. While pretraining and post-training are functionally linked, current practices often treat them as decoupled. Specifically, pretraining hyperparameters and scaling laws are predominantly studied through the lens of the base model's validation loss, under the assumption that a lower pretraining validation loss also yields a more capable downstream model \citep{hoffmann2022training,bi2024deepseek}. This decoupling is especially pronounced when the two stages are carried out by different teams or at different times: a ``best'' pretrained model is selected in isolation and only later adapted for downstream use. However, to what extent does optimizing pretraining hyperparameters for pretraining performance also optimize the final, post-trained model's performance?

In this work, we study the relationship between pretraining and post-training from the perspective of \textit{model plasticity}. Model plasticity is the ability of a trained model to effectively adapt to new data upon further training, modifying its parameters and internal representations in response to the new data and enabling effective learning of new tasks without reinitialization~\citep{berariu2021study,dohare2024loss}. 
While the literature on model plasticity and on language models have evolved largely independently, the notion of plasticity naturally bridges pretraining and post-training: while pretraining loss measures how well a model learned the training distribution, plasticity captures how readily that model can be reshaped for downstream tasks. As we show in this work, two models with similar pretraining loss may differ in their plasticity, meaning that optimizing hyperparameters for pretraining loss alone may not yield the best post-trained model.


In our experiments, we vary the weight decay hyperparameter during pretraining and subsequently evaluate the pretrained model’s ability to learn various tasks during fine-tuning. Pretraining is performed for two model families (Llama-2 and OLMo-2), multiple model sizes (up to 4B parameters), and in both the compute-optimal (20 tokens-per-parameter, TPP hereafter) and overtrained (140 TPP) regimes. Fine-tuning is performed across six Chain-of-Thought (CoT) reasoning tasks, five language understanding and commonsense reasoning tasks, and one safety alignment task, and model performance is evaluated using a comprehensive suite of metrics that cover both solution correctness and quality. Our experimental design takes an end-to-end perspective \citep{qi2025evolm,mayilvahanan2025llmsontheline}, aligning pretraining hyperparameter selection with the ultimate objective of maximizing performance after further training.
Our contributions are as follows:

\begin{itemize} 
\item We show that weight decay is a key factor in shaping model plasticity, facilitating adaptation to new tasks during subsequent fine-tuning. In our experiments across a range of model families, sizes, training regimes, downtream tasks, and evaluation metrics, the evidence points toward an optimal pretraining weight decay value larger than the standard default of 0.1. This highlights the potential for re-evaluating standard hyperparameter choices to better account for a model's downstream adaptability.

\item We provide one of the first examples showing that optimizing hyperparameters to minimize pretraining validation loss does not necessarily yield the best downstream model performance. Specifically, we show that there is a training regime where larger weight decay values lead to higher pretraining validation loss {\it and} better downstream performance after fine-tuning.

\item We provide a mechanistic perspective on the effect of weight decay on model training dynamics, showing that weight decay encourages linearly separable representations, regularizes attention matrices, and reduces overfitting on the training data. These effects provide a potential explanation for how weight decay preserves the model's ability to flex and learn during subsequent adaptation, thereby sustaining plasticity and improving downstream performance. 
\end{itemize}


Together, these findings highlight the importance of pretrained model plasticity in the pretrain-then-post-train development pipeline of modern language models, the limits of using cross-entropy loss as the sole metric for hyperparameter optimization during pretraining, and the multifaceted role of a single optimization hyperparameter (weight decay) in shaping model behavior.